\chapter{}

Pietro si commoveva fino a pensare: «Se anche fosse disonesta per necessità di non patire la fame, io non potrei approfittarne. Piangerei. L'aiuterei a fare in modo che si cambiasse. Qualcuno, allora, potrebbe stimarla e sposarla. Ma me lo avrebbe detto. Perchè non me lo dovrebbe dire?».

E, per contrasto al dubbio, gli pareva d'una purità mirabile. Allora ne era geloso e piangeva. «Deve esser mia! Voglio amarla io! Perchè non dovrei amarla?» Non era anche il suo dovere morale? Ma come il modo di star meglio che in casa del padre? Ghìsola gli aveva detto:

— È ricco; dipende tutto da lui. Ma egli non vorrà di certo.

Domenico, quando Pietro, tornato da Firenze, gli disse ch'era innamorato di Ghìsola e che, se fosse stato contento, aveva deciso di sposarla, non gli rispose nè meno; ma si sentì aizzato contro di lui come la volpe quando le hanno accesa la paglia dentro la tana.

Degli esami tacquero ambedue. Pietro per non fargli sapere la verità, e Domenico per tentare che non ci pensasse più, ma con la voglia di sbatterlo nel muro come un cuscino.
